\exercises

\tierunderstand

\begin{bt}
  \item Dẫn \eqref{eq:ch08-flop-tang} lại từ đầu, viết rõ hình dạng của từng ma
    trận ở từng phép nhân. Rồi trả lời câu đắt hơn: nếu bốn phép chiếu
    \emph{có} bias thì số hạng nào đổi, và đổi bao nhiêu? Con số ấy có làm
    \eqref{eq:ch08-giao-diem} dịch đi không?

  \item Mục~\ref{sec:ch08-dem} nói số head triệt tiêu khỏi phép đếm FLOP. Chỉ ra
    đúng bước đại số nơi nó triệt tiêu, và nói bước ấy dựa vào giả thiết nào.
    Rồi bỏ giả thiết đi: nếu một cài đặt đặt \(\dk\) là một hằng số thay vì
    \(\dmodel/h\), công thức thành gì, và \code{tests/test\_ch08.py} sẽ hỏng ở
    test nào?

  \item Từ \eqref{eq:ch08-ty-le}, chứng minh tỷ lệ bậc hai chỉ phụ thuộc
    \(n/(2d + \dff)\). Rồi dùng nó trả lời mà không chạy gì: ở
    \(\dmodel = 1024\) và \(\dff = 4096\), phải lấy \(n\) bằng bao nhiêu thì
    phần bậc hai chiếm 90\% một tầng?

  \item Mục~\ref{sec:ch08-dem} đo được attention rẻ hơn LSTM tới \(n = 2d\) mà
    lại đắt hơn một tầng hồi quy trơn ở mọi \(n\). Giải thích cả hai bằng
    \eqref{eq:ch08-flop-tang}. Rồi nói bảng 1 của bài lẽ ra phải viết thế nào
    để hai chuyện ấy đọc ra được từ bảng.

  \item Mục~\ref{sec:ch08-cho-vo} đo được cái dốc đi theo số byte chứ không theo
    \(n\). Thiết kế một phép đo thứ ba tách hai giả thuyết ấy, khác cả phép quét
    \(n\) lẫn phép quét batch mà mục ấy đã chạy. Nói trước phép đo của bạn sẽ
    cho kết quả gì dưới mỗi giả thuyết.

  \item Ở \eqref{eq:ch08-warmup}, chứng minh hai nhánh của \(\min\) bằng nhau
    đúng tại \(s = w\), và tính giá trị đỉnh theo \(\dmodel\) và \(w\). Rồi trả
    lời: muốn giữ nguyên learning rate đỉnh khi tăng \(\dmodel\) gấp bốn, phải
    đổi \(w\) thế nào?

  \item Định lý 1 ở hộp cầu nối mục~\ref{sec:ch08-thu-tu-ln} là một upper bound.
    Giải thích vì sao một phép đo cho thấy gradient pre-LN \emph{không} giảm
    theo \(1/\sqrt{L}\) thì không bác được định lý. Rồi mô tả một kết quả đo
    nào \emph{sẽ} bác được nó.

  \item Mô hình pre-LN nhiều hơn post-LN đúng \(4\dmodel\) tham số. Chỉ ra bốn
    nhóm ấy nằm ở đâu. Rồi nói vì sao một phép so hai kiến trúc chỉ nhìn số tham
    số sẽ coi hai mô hình này là một, và liên hệ với phép cắt bỏ có đối chứng
    cùng số tham số ở chương~\ref{ch:mot-vector-cho-moi-tu}.

  \item \eqref{eq:ch08-chan-duoi} cho lower bound của mất mát khi có label
    smoothing. Kiểm nó ở một trường hợp bạn tính tay được: \(K = 2\),
    \(\epsilon = 0.1\). Rồi nói lower bound ấy đi về đâu khi \(K\) lớn dần với
    \(\epsilon\) giữ nguyên, và chuyện đó
    nói gì về việc so mất mát giữa hai mô hình khác kích thước từ điển.
\end{bt}

\tierapply

Trạng thái xuất phát cho cả năm bài: clone repo companion, đứng ở thư mục gốc
của nó, và checkout đúng tag của chương này.

\begin{minted}{console}
$ git clone https://github.com/Giang-Dang/rnn-to-transformer-lab.git
$ cd rnn-to-transformer-lab
$ git checkout ch08
$ conda env create -f environment.yml
$ conda activate rnn-to-transformer-lab
$ python verify.py --only ch08
\end{minted}

Dòng cuối phải in ra \code{verify: ok}. Cả repo chạy mất khoảng hai mươi phút;
\code{--only ch08} chạy riêng năm mục của chương này và mất sáu tới bảy. Hai con
số ấy là của máy tôi đo, và chúng dao động theo việc máy có đang chạy gì khác
không: hai lần chạy cả repo trên cùng mã trong cùng một buổi cách nhau hơn hai
phút.

Vẫn cảnh báo cũ: đừng gọi các script qua \code{conda run}, vì nó gom stdout của
tiến trình con rồi in lại bằng encoding hệ thống và ném
\code{UnicodeEncodeError} từ bên trong chính conda.

\begin{bt}
  \item Đổi quy ước đếm. \code{cost.py} đếm một phép nhân-cộng là hai FLOP. Sửa
    thành một, chạy lại \code{ch08\_flops.py}, và nói bảng nào đổi và bảng nào
    không. Rồi chạy \code{python -m pytest -q tests/test\_ch08.py} và giải thích
    vì sao đúng một test hỏng chứ không phải bốn.

  \item Đo lại cái dốc trên máy của bạn. Chạy \code{ch08\_clock.py} và tìm chỗ
    cột \code{s/GFLOP} bắt đầu leo. Tra kích thước L3 của CPU bạn dùng
    (\code{wmic cpu get L3CacheSize} trên Windows, \code{lscpu} trên Linux) rồi
    đối chiếu với cột \code{score MiB}. Chỗ dốc của bạn có rơi vào đúng chỗ
    cache của bạn hết chỗ không? Nếu không, nêu một cách giải thích khác và một
    phép đo tách nó.

  \item Cho pre-LN đủ độ sâu. \code{ch08\_norm.py} huấn luyện ở 2 tầng, và
    mục~\ref{sec:ch08-warmup} nói đó là lý do bảng bốn ô không kết luận được.
    Sửa cho nó chạy \(L\) bằng 2, 4 và 8 ở cả hai thứ tự, ba seed mỗi ô, rồi hỏi
    khoảng \code{min}-\code{max} có tách ra ở \(L\) nào không. Bài này vượt xa
    ngân sách của \code{verify.py}; chạy ngoài nó và nói rõ tổng thời gian bạn
    đã tiêu.

  \item Bỏ cái layer norm cuối cùng của pre-LN. Trong \code{transformer.py},
    đổi \code{encoder\_norm} và \code{decoder\_norm} thành \code{nn.Identity()}
    ở cả hai nhánh, rồi chạy \code{tests/test\_ch08.py} và
    \code{ch08\_norm.py}. Test nào bắt được? Bảng gradient lúc khởi tạo đổi thế
    nào, và mô hình huấn luyện ra sao? Trả lời câu cuối: nếu chỉ được nhìn
    đường cong mất mát, bạn có phát hiện ra lỗi này không?

  \item Đẩy label smoothing tới chỗ nó hỏng. \code{ch08\_recipe.py} chạy
    \(\epsilon_{ls}\) bằng 0.0 và 0.1. Thêm 0.3, 0.5 và 0.9, giữ nguyên ba seed,
    và in cả cột \code{excess}. Có giá trị nào mà khớp đúng cả câu tụt hẳn
    không? Đối chiếu chỗ ấy với hàng \(\epsilon_{ls} = 0.2\) của bảng 3, nơi bài
    đo BLEU 25.7 so với 25.8 của base, và nói vì sao hai bảng không so trực tiếp
    được với nhau.
\end{bt}

\tierextend

\begin{bt}
  \item Mục~\ref{sec:ch08-dong-ho} in hai cột cho hai cách đọc chữ
    \enquote{recurrent}, và vòng lặp Python của repo này là cột chậm. Câu hỏi:
    bao nhiêu phần của khoảng cách ấy là bản chất của tính tuần tự, và bao nhiêu
    là chi phí của trình thông dịch? Thiết kế một phép đo tách hai phần, ví dụ
    bằng cách viết cùng một vòng lặp ấy dưới \code{torch.compile} hoặc
    \code{torch.jit.script} và so ba cột thay vì hai. Kết luận của bạn đổi cách
    đọc bảng 1 chỗ nào?

  \item Đây là chỗ còn mở thật sự, và mục~\ref{sec:ch08-thu-tu-ln} dừng lại ở
    đúng mép của nó. Định lý 1 chứng minh trong cấu hình đặt
    \(\mat{W}^Q = \mat{W}^K = 0\), tức attention là trung bình cộng đều ở lúc
    khởi tạo. Câu hỏi: upper bound ấy còn đúng dạng nào khi \(\mat{W}^Q\) và
    \(\mat{W}^K\) khởi tạo như \code{nn.Linear} thật sự làm, tức đều trên
    \([-1/\sqrt{n_{\mathrm{in}}}, 1/\sqrt{n_{\mathrm{in}}}]\)? Làm được một
    trong hai chiều là đủ: dẫn lại upper bound ấy với giả thiết khởi tạo mới,
    hoặc dựng một dãy mô hình mà số đo tách hẳn khỏi \(1/\sqrt{L}\) đủ để nói
    upper bound cũ không mô tả chúng. Biết trước: mục~\ref{sec:ch07-chia-can} đã đo rằng khởi tạo ấy cho
    phương sai thành phần 1/3 chứ không phải 1, nên hằng số sẽ không sạch.

  \item Bảng 1 của bài có ba cột, tất cả đều là hàm của \(n\), \(d\), \(k\) và
    \(r\). Mục~\ref{sec:ch08-cho-vo} đo được một đại lượng không viết được thành
    hàm của bốn tham số ấy, vì nó phụ thuộc cache của máy. Đề xuất một cột thứ
    tư cho bảng 1 mô tả chi phí truy cập bộ nhớ, định nghĩa nó đủ chặt để tính
    được cho cả bốn hàng, rồi tính. Nếu cột của bạn phải mang thêm một tham số
    của phần cứng, hãy nói rõ đó là tham số nào, và trả lời câu thật sự khó:
    một bảng \tn{độ phức tạp}{complexity} mang tham số phần cứng thì còn là một
    bảng độ phức tạp
    nữa không?

  \item Mục~\ref{sec:ch08-dieu-chuan} đo được dropout làm hỏng ở đây và bài đo
    được nó đáng 1.2 BLEU. Chương này quy chỗ khác nhau cho lượng dữ liệu, và
    đó là một giả thuyết chưa được kiểm ở đây. Kiểm nó: giữ nguyên mọi thứ và
    quét số cặp huấn luyện từ 6000 lên tới chỗ máy bạn còn chịu được, ở cả
    \(P_{drop} = 0.0\) và \(0.1\), rồi hỏi hai đường có cắt nhau không. Nếu
    không cắt, giả thuyết \enquote{tại ít dữ liệu} sai, và lời giải thích còn
    lại là văn phạm hữu hạn sinh ra corpus chứ không phải số lượng; thiết kế
    tiếp một phép đo tách hai thứ ấy.
\end{bt}
